\section{Further research questions and a prioritized program}
\label{sec:agenda}

The most useful next step is not simply a larger benchmark. It is to
separate structural hypotheses that might yield a quasi-polynomial theorem
from engineering changes whose value must be established by controlled
whole-query experiments. The questions below specify the desired object,
the proof obligation, and a useful first deliverable.

\subsection{Structural questions bearing directly on the target}

\begin{enumerate}[label=\textbf{Q\arabic*.},leftmargin=3em]
\item \textbf{Which precisely defined residual class admits small
presentation parameters?}
After certified prime decomposition and the maintained structural,
determinant, coloring, and Alexander tests, can the surviving class be
shown to have a constructible circuit cut satisfying
$(r+t)\log(d+2)=O(\log^2 n)$? The filters must be fixed mathematically,
not described as ``all easy cases.'' Begin with prime determinant-one or
Alexander-trivial diagrams and identify whether large ordinary group
rank persists. The displayed dihedral obstructions are removed by cheap
coloring and therefore do not settle this residual question. A counterfamily
would be as useful as a positive theorem, provided it survives the stated
filters and includes a rank or cut lower bound.

\item \textbf{Can one construct small presentations of unknots efficiently?}
Every unknot group has a one-generator presentation, but that existence
does not bound the cost of finding it from a difficult diagram. Seek a
family of verified Tietze or geometric transformations with bounded move
count, grammar growth, and compressed verification cost on unknots.
Separate unknotted-input guarantees from the separate task of rejecting
knotted inputs. A proof that a local search never stalls requires a
global completeness argument, and any temporary word expansion must be
charged in its actual representation.

\item \textbf{Can the decision problem use smaller data than the whole group?}
The rank obstructions rule out a universal small presentation of the
entire group, not a small certificate tailored to triviality. One could
seek verified decompositions, representation-preserving quotients, or
local extension problems whose conjunction detects the unknot. A proposed
quotient must preserve the direction needed: a nonabelian image of a
quotient certifies knottedness, whereas failure to find one generally
does not certify the unknot. A complete reduction needs both directions
and a bound on the size and number of subproblems.

\item \textbf{Can checkpoint gaps and frontier widths be controlled together?}
Construct crossing orders or certified splicing moves with
$g+w=O(\log^2 n)$ on a meaningful class of hard residual diagrams.
The gap must include the crossing that triggers a replacement and the
last segment; the width must be the actual temporary frontier. Study
whether alternating reductions, arborescent decompositions, or controlled
tangle sums give such orders after the already easy cases are removed.
A bound only on final survivor count leaves the intermediate cost unresolved.

\item \textbf{Can geometric hierarchy certificates supply the missing size theorem?}
Translate the hierarchy parameters in \cite{Lackenby2026Hierarchies} into
explicit data structures and verified transitions. Bound hierarchy length,
surface and pattern complexity, compressed cutting cost, and the bit
length of every update. The displayed iteration bound $L(g+1)^L$ is
only one term in that analysis. A practical starting point is a small
certificate checker for one cutting or simplification step, with
independent topology reconstruction and a precise input-size theorem.

\item \textbf{What exactly is established by the recent bridge-minimization claim?}
Independently audit \cite{Musick2026}, including conversion from arbitrary
input diagrams, completeness of the allowed moves, the application of
Otal's theorem to the unknot case, and polynomial
bounds for their actual encoded
execution. Record any lemma that requires an additional hypothesis and
distinguish implementation tests from proof of that lemma. The present
report's results do not assume the claim; resolving its status would
materially change the broader research landscape.
\end{enumerate}

\subsection{Algebraic compilation and exact certificates}

\begin{enumerate}[label=\textbf{Q\arabic*.},leftmargin=3em,start=7]
\item \textbf{How hard is the best word-circuit cut?}
Study the optimization problem
$\min_C(r+|C|)\log(d_C+2)$ on a word DAG. A first target is an exact
dynamic program for trees, followed by a proved approximation or a
parameterized algorithm for shared DAGs. Repeated subwords make local
greedy cuts potentially suboptimal. Formal degree, actual monomial
support, coefficient height, and compiler work can be tracked separately;
an empirical support heuristic should not be substituted into the
worst-case theorem without a bound.

\item \textbf{Can quaternion-specific identities improve the degree bound?}
Repeated powers of a unit quaternion satisfy a quadratic recurrence, so
Chebyshev-type representations may compress certain relators further.
Any use of these identities must preserve all unit-quaternion solutions,
including central generators, and count newly introduced scalar or
quaternion variables. Determine whether power-heavy knot presentations
admit smaller cuts than generic straight-line words. Avoid confusing a
short recurrence with a constant-cost real formula.

\item \textbf{Can exact representation witnesses be checked independently?}
The current optional solver records algebraic models and checks its
assertions through the same exact engine. A stronger SAT artifact would
give isolating intervals, defining polynomials, and independently verified
sign tests for every equation and the noncommutativity inequality. For
UNSAT, specify a checkable real-algebraic proof format or a verified
decision backend. Establish verifier complexity in the certificate's
actual bit size; a compact approximate numerical model alone is insufficient.

\item \textbf{How should the known traceless-meridian theorem be used?}
Kronheimer--Mrowka already prove complete detection with a meridian
sent to $i$ \cite[Corollary~7.17]{KM2010Excision}; this is not an open
trace-existence question. Implement and independently validate
Corollary~\ref{cor:su2-fixed-meridian}, carrying explicit meridional
provenance through generator changes. Compare its smaller complete formula
with the current one-sided witness probe on the same residual corpus.
A separate mathematical question concerns other prescribed trace values
or trace schedules that improve real-algebraic decomposition. Such a
restriction requires its own theorem, and generic nonmeridional
generators must retain the unrestricted query.

\item \textbf{Can presentation provenance remain compressed across every move?}
The repository already independently replays several compressed group
certificates. Connect that verifier to the new circuit compiler so arbitrary
generator changes switch correctly from the meridional inequality to the
general commutator condition. Bound grammar size and assertion growth
through long traces, not just individual word operations. The first
deliverable should be a shared, versioned presentation object carrying
origin, generator meanings, and exactly the current relators.
\end{enumerate}

\subsection{Earlier homological decisions and reusable geometry}

\begin{enumerate}[label=\textbf{Q\arabic*.},leftmargin=3em,start=12]
\item \textbf{Can the flat-origin constraint remove the second rank safely?}
For genuine minimal same-matching components, $\kappa=t$. Determine how
to maintain a cheap connection or provenance certificate across scalar
cancellation, direct-summand extraction, and total-rank replacements.
The last operation need not preserve genuine cube origin. A theorem
covering exactly the executed scanner would permit using one rank where
two are currently computed. Merely observing $t=\kappa$ in a corpus is
not sufficient to install that shortcut.

\item \textbf{What constraints survive in higher jets or mixed matchings?}
Lift the identity $\partial Q=\Gamma Q+Q\Gamma$ beyond scalar residue.
Can it restrict higher-dot coefficients, transferred operations, or the
shape of mixed-matching minimal blocks enough to bound multiplicities?
The objective is a genuine realizability constraint on tangle-derived
complexes. Arbitrary polynomial complexes over the dot algebra may be
too broad to support the desired theorem, as the synthetic odd intervals
already demonstrate.

\item \textbf{Can eligible components be maintained incrementally?}
Current graph traversal and binary ranks can cost more than the scalar
cancellations they avoid. Develop a dynamic index that certifies a
whole one-matching direct summand and updates its scalar matrix after one
crossing. It must detect all external attachments and survive transactional
interruption. Evaluate saved cobordism compositions against the cost of
component maintenance; the present easy replacement families regress.

\item \textbf{When should a scan precompute suffix responses?}
Develop an adaptive policy that chooses independent setup, checkpoint
responses, or full reverse compilation based on a proved cost estimate
and observed demand. Aim for a competitive bound relative to the better
strategy, including source geometry and abandoned work. A forward scanner
may terminate before most precomputed suffixes are queried. The large
all-stage timing ratios therefore cannot directly choose a default policy.

\item \textbf{Can a separator tree support local response updates?}
Replace the linear order by a decomposition tree while retaining radical
couplings and exact grounding. Prove a boundary contract that permits
subtree changes without identifying eliminated interior vertices. Determine
which snapshots and Euler summaries can be shared across primes or nearby
orders despite exceptional pivot profiles. The established low-treewidth
linear-algebra literature provides a starting point, not a substitute for
the topology-specific gluing proof \cite{furer-hoppen-trevisan}.
\end{enumerate}

\subsection{Evidence, dispatch, and a concrete sequence of deliverables}

\begin{enumerate}[label=\textbf{Q\arabic*.},leftmargin=3em,start=17]
\item \textbf{Which residual corpus measures the intended improvement?}
Assemble difficult unknots and knots that survive the exact existing
front ends, with full PD provenance, mirrors, several orders, and
duplicate detection. Include determinant-one and Alexander-trivial knots,
inflated unknots, and examples with substantial repeated suffix interior
work. A test label should come from an independent certificate or complete
oracle. Synthetic word circuits remain useful for capacity tests but
must be reported separately from knot-recognition instances.

\item \textbf{Which optional method improves complete queries reliably?}
Run paired, shuffled whole-recognizer experiments with identical controls,
fixed source hashes, explicit warmup rules, and all outcomes retained.
Report CPU and wall time scopes, object and coefficient peaks, solver
startup, model validation, and censored runs. Promotion to a default
should require a reproducible benefit on the target residual corpus
without a correctness or deadline regression. A small local improvement
does not justify hiding a larger all-stage setup cost.
\end{enumerate}

The proposed order is concrete. First, finish the residual corpus and
attach independently checked presentation provenance to the representation
compiler. Next, reduce first-jet traversal overhead and prototype an
adaptive response consumer, retaining the present backends as controls.
In parallel, pursue either a circuit-cut theorem for a precisely filtered
class or a geometric bound on checkpoint gaps and frontier width. A
successful structural theorem would immediately plug into
Corollary~\ref{cor:regimes}; until then, the optional implementations and
recorded negative results provide a disciplined experimental platform.
